follows from the uniqueness of integer division and congruence class representatives. The orbits under the translation group action coincide with congruence classes, $\mathcal F_b$ is a transversal thereof, and the projection $\Pi_b$ yields a unique $r$, from which $q$ is determined.
\end{proof}

\begin{corollary}[Centered remainder transversal]\label{cor:pom-symmetric-remainder}
If one instead uses the symmetric fundamental domain
\[
\mathcal F_b^{\mathrm{sym}}:=\left\{-\left\lfloor\frac{b}{2}\right\rfloor,\dots,\left\lceil\frac{b}{2}\right\rceil-1\right\},
\]
then the unique decomposition $a=qb+r$ with $r\in\mathcal F_b^{\mathrm{sym}}$ still holds. This choice amounts to changing the transversal rather than the ontological dynamics, and can consolidate part of the ``interval correction'' into a single order projection.
\end{corollary}

\begin{remark}[Groupoid analogy with Fold]
Theorem \ref{thm:pom-division-section} and Theorem \ref{thm:zeckendorf-transversal} are structurally isomorphic: both exhibit the mechanism ``orbit = value fiber / congruence class, transversal = canonical representative''; the quotient $q$ and the groupoid coordinate of the folding residual occupy the same semantic position.
\end{remark}

\subsection{Conclusion Cluster III: Extended Minimal Primitive Set Including Quotient and Remainder}\label{subsec:pom-primitive-extended}
\begin{theorem}[Extended minimal primitive set]\label{thm:pom-extended-primitives}
At the configuration level, $\oplus,\otimes,\mathrm{div},\mathrm{quo},\mathrm{rem}$ can be generated by the following primitives:
\begin{enumerate}
  \item Pointwise superposition (free commutative monoid structure);
  \item Fibonacci groupoid generators (value-preserving locally invertible rewrites);
  \item Value fiber transversal projection $\Fold_\infty$ (canonical projection);
  \item Translation groupoid transversal projection $\Pi_b$ (order/interval projection).
\end{enumerate}
At the process level, the equivalent formulation is: using only composition $\circ$ as the underlying primitive, augmented by the value projection $\Fold_\infty$ and the order projection $\Pi_b$ to land the output on the external surface; the atomic projection $P_{\mathrm{prim}}$ is applied independently as a spectral readout interface.
\end{theorem}
\begin{proof}
Addition and multiplication are realized at the process level by the readout of Theorem \ref{thm:composition-two-layer} and \ref{thm:arith-composition}, and are landed on stable addresses by $\Fold_\infty$. Quotient and remainder are provided by Theorem \ref{thm:pom-division-section} as the translation groupoid transversal projection. The remaining operations are obtained from the composition closure of the above primitives.
\end{proof}

\begin{theorem}[Double transversal normal form: value fiber $\times$ congruence class]\label{thm:pom-double-transversal-normal-form}
For any configuration object $a\in\mathcal D$ and any integer $b>0$, there exists a unique triple
\[
\boxed{(x,q,r)\in X_\infty\times\mathbb{Z}\times\mathcal F_b}
\]
such that
\[
x=\Fold_\infty(a),\qquad
r=\Pi_b(\mathrm{Val}(a)),\qquad
q=\frac{\mathrm{Val}(a)-r}{b},
\]
and \(\mathrm{Val}(a)=\mathrm{Val}(x)=qb+r\).
\end{theorem}
\begin{proof}
By Theorem \ref{thm:zeckendorf-transversal} (or equivalently the definition \(\Fold_\infty(a)=Z(\mathrm{Val}(a))\)), \(x\) is unique and satisfies \(\mathrm{Val}(x)=\mathrm{Val}(a)\).
By Theorem \ref{thm:pom-division-section}, \((q,r)\) is unique and satisfies \(r=\Pi_b(\mathrm{Val}(a))\) and \(\mathrm{Val}(a)=qb+r\).
Combining both parts yields the existence and uniqueness of the stated triple.
\end{proof}

\begin{corollary}[Constant budget for projection invocations: terminal normalization]\label{cor:pom-projection-budget}
Let \(t\) be any expression composed from configuration-level operations \(\oplus,\otimes,\mathrm{div},\mathrm{quo},\mathrm{rem}\) (of arbitrary depth), and suppose its \emph{externally visible output} consists of a single stable address readout (landing in \(X_\infty\)) together with an optional quotient-remainder readout (landing in \(\mathbb Z\times\mathcal F_b\)).
Then there exists an equivalent process-level implementation in which, throughout the entire evaluation chain:
\begin{itemize}
  \item The value transversal projection \(\Fold_\infty\) is invoked at most once (only at the terminal step to land the output in \(X_\infty\));
  \item The order transversal projection \(\Pi_b\) is invoked at most once (only when the output type requires a unique remainder representative, applied at the terminal step to \(\mathrm{Val}\)).
\end{itemize}
Moreover: if the output type requires a unique stable address (or a unique remainder representative), then the corresponding terminal projection is indispensable in that sense.
\end{corollary}
\begin{proof}
By Theorem \ref{thm:pom-extended-primitives}, the process-level evaluation can be rewritten as the composition of invertible evolution within the groupoid and terminal transversal projections.
Since \(\Fold_\infty\) is the canonical section of \(\mathrm{Val}\): \(\mathrm{Val}(\Fold_\infty(a))=\mathrm{Val}(a)\), and \(\Fold_\infty(a)\) depends only on \(\mathrm{Val}(a)\),
any intermediate normalization that repeatedly ``lands and resumes evolution'' within the same value fiber does not alter the terminal output that lands in \(X_\infty\) according to the final value, and can therefore be uniformly deferred to a single terminal invocation.
Similarly, \(\Pi_b\) depends only on the integer value \(\mathrm{Val}(\cdot)\), and \(\mathrm{Val}\) can be read out once at the terminal step; hence if only a single quotient-remainder readout is needed, \(\Pi_b\) can likewise be deferred to a single terminal invocation.
Finally, the ``indispensability'' follows from the non-uniqueness of orbit representatives: without applying the corresponding transversal projection, the output remains on the orbit rather than at a unique point.
\end{proof}

\subsection{Inverse Limit Naturality and Fiber Invariant Discrimination}\label{subsec:pom-new-results}
\subsubsection{The Gauge Principle for Cross-Resolution Consistency: Fold-Aware Restriction}\label{subsubsec:pom-gauge}
The finite-window microstate space $\Omega_m=\{0,1\}^m$ and the stable type space $X_m$ play asymmetric roles: $X_m$ forms an inverse system under the truncation maps $\pi_{m+1\to m}$, whereas $\Omega_m$ does not carry an inverse system structure compatible with folding under naive truncation. This asymmetry is precisely the structural origin of the requirement that ``cross-resolution consistency must first be normalized.''

\begin{definition}[Naive truncation and fold-aware restriction]\label{def:pom-fold-aware-restriction}
Let $\tau_{m+1\to m}:\Omega_{m+1}\to\Omega_m$ be the naive truncation (discarding the last bit; equivalently taking the length-$m$ prefix). Define the fold-aware restriction as
$$
r^{\mathrm{e}}_{m+1\to m}:\Omega_{m+1}\to X_m,
\qquad
r^{\mathrm{e}}_{m+1\to m}:=\pi_{m+1\to m}\circ\Fold_{m+1}.
$$
\end{definition}

\begin{proposition}[Naive truncation and folding generally do not commute]\label{prop:pom-truncation-not-commute}
In general, the diagrams
$$
\Omega_{m+1}\xrightarrow{\ \tau_{m+1\to m}\ }\Omega_m\xrightarrow{\ \Fold_m\ }X_m
$$
and
$$
\Omega_{m+1}\xrightarrow{\ \Fold_{m+1}\ }X_{m+1}\xrightarrow{\ \pi_{m+1\to m}\ }X_m
$$
do not agree, i.e., there exists $\omega\in\Omega_{m+1}$ such that
$$
\Fold_m(\tau_{m+1\to m}(\omega))\neq \pi_{m+1\to m}(\Fold_{m+1}(\omega)).
$$
\end{proposition}
\begin{proof}
The local rewriting of the fold (e.g., $011\mapsto 100$) produces cross-bit propagation: the discarded last bit may determine whether a particular local rewrite is triggered, thereby altering the normalized form of the length-$m$ prefix. Consequently, ``truncate first, then fold'' and ``fold first, then truncate'' are generally inconsistent.
\end{proof}

\begin{proposition}[Naturality of the stable inverse system: fold first, then restrict]\label{prop:pom-gauge-natural}
For any $m\ge 1$, the fold-aware restriction satisfies
$$
r^{\mathrm{e}}_{m+1\to m}
=
\pi_{m+1\to m}\circ \Fold_{m+1},
$$
so that at the stable level the following cross-resolution diagrams strictly commute:
$$
\Omega_{m+1}\xrightarrow{\ \Fold_{m+1}\ }X_{m+1}\xrightarrow{\ \pi_{m+1\to m}\ }X_m
\qquad\text{and}\qquad
\Omega_{m+1}\xrightarrow{\ r^{\mathrm{e}}_{m+1\to m}\ }X_m
$$
are identical; and for any $m_2\ge m_1$, the compatibility
$$
\text{let }r^{\mathrm{e}}_{m_2\to m_1}:=\pi_{m_2\to m_1}\circ\Fold_{m_2},
\ \text{then for any }m_3\ge m_2\ge m_1\text{ we have }r^{\mathrm{e}}_{m_3\to m_1}=\pi_{m_2\to m_1}\circ r^{\mathrm{e}}_{m_3\to m_2},
$$
holds, so that ``cross-resolution coordinates of a single object'' can only be discussed within the inverse system $X_\infty\simeq\varprojlim X_m$.
\end{proposition}
\begin{proof}
The first identity is simply the expansion of Definition \ref{def:pom-fold-aware-restriction}. Compatibility follows from the universal property satisfied by the inverse system maps at the stable level (Theorem \ref{thm:pom-inverse-limit}), and the coordinates after folding are precisely the inverse system objects.
\end{proof}

\begin{definition}[Gauge curvature (2-cocycle obstruction)]\label{def:pom-gauge-curvature}
For a given resolution descent \((m{+}1\to m)\) and a top-level microstate \(\omega\in\Omega_{m+1}\), define the curvature indicator
\[
\mathcal{K}_{m+1\to m}(\omega)
:=
\mathbf 1\!\left[\,
\Fold_m(\tau_{m+1\to m}(\omega))
\neq
\pi_{m+1\to m}(\Fold_{m+1}(\omega))
\,\right].
\]
It quantifies the non-commutativity obstruction between ``truncate first, then fold'' and ``fold first, then restrict'' at this microstate.
\end{definition}

\begin{proposition}[Measurable error band for pseudo-invariants: controlled by curvature probability]\label{prop:pom-gauge-curvature-shadow-bound}
Let \(\mu\) be any distribution on \(\Omega_{m+1}\), and let \(f:X_m\to\mathbb{R}\) be any bounded observable.
Then the following auditable inequality holds:
\[
\boxed{
\begin{aligned}
&\left|\mathbb E_{\omega\sim\mu}\,f(\Fold_m(\tau_{m+1\to m}(\omega)))
-\mathbb E_{\omega\sim\mu}\,f(\pi_{m+1\to m}(\Fold_{m+1}(\omega)))\right|\\
&\qquad \le 2\|f\|_\infty\cdot \mathbb P_{\omega\sim\mu}\!\left[\mathcal{K}_{m+1\to m}(\omega)=1\right].
\end{aligned}}
\]
In particular, under the gauge convention of fold-aware restriction (the strictly commuting diagram of Proposition \ref{prop:pom-gauge-natural}), the above error band collapses to \(0\).
\end{proposition}
\begin{proof}
For each \(\omega\), let
\(
x=\Fold_m(\tau_{m+1\to m}(\omega))
\)
and
\(
y=\pi_{m+1\to m}(\Fold_{m+1}(\omega))
\).
Then pointwise,
\[
|f(x)-f(y)|
\le
2\|f\|_\infty\,\mathbf 1[x\neq y]
=
2\|f\|_\infty\,\mathcal{K}_{m+1\to m}(\omega).
\]
Taking the expectation over \(\omega\sim\mu\) yields the stated inequality; the final assertion follows directly from Proposition \ref{prop:pom-gauge-natural}.
\end{proof}

\begin{remark}[Reproducible audit: enumeration/sampling of curvature probability]\label{rem:pom-gauge-curvature-audit}
The curvature probability on the right-hand side,
\(\mathbb P[\mathcal{K}_{m+1\to m}=1]\),
can be directly obtained by simultaneously computing the outputs of both paths and tallying mismatch frequencies.
The script \path{scripts/exp_fold_truncation_curvature.py} performs exact enumeration under the uniform microstate baseline for \(m\le 20\) and generates the table fragment; for general data sources the same definition can be estimated by sampling.
For instance, the counting functions used in the main text to define prime shadows, or any cylinder indicator function \(f\), are bounded observables on stable types, and hence the inequality of Proposition \ref{prop:pom-gauge-curvature-shadow-bound} provides an auditable ``artifact error band.''
\end{remark}

\input{sections/generated/tab_fold_truncation_curvature_stats_en}

\begin{definition}[Resolution reflector (cross-resolution fiber uniformization)]\label{def:pom-resolution-reflector}
Fix a top-level resolution $M$ and take the microstate space as $\Omega_M$.
For any $1\le m\le M$, using the fold-aware restriction
\(
r^{\mathrm{e}}_{M\to m}=\pi_{M\to m}\circ \Fold_M:\Omega_M\to X_m
\)
as the coarse-graining map, define the pushforward operator
\[
P_{M\to m}:=(r^{\mathrm{e}}_{M\to m})_*:\Delta(\Omega_M)\to\Delta(X_m),
\]
and denote its fiber cardinality by
\[
d_{M\to m}(x):=\bigl|(r^{\mathrm{e}}_{M\to m})^{-1}(x)\bigr|\qquad(x\in X_m).
\]
Define the fiber-uniform lift operator
\[
Q_{M\to m}:\Delta(X_m)\to\Delta(\Omega_M),
\qquad
(Q_{M\to m}\pi)(\omega):=\frac{\pi(r^{\mathrm{e}}_{M\to m}(\omega))}{d_{M\to m}(r^{\mathrm{e}}_{M\to m}(\omega))}.
\]
Then $P_{M\to m}\circ Q_{M\to m}=\mathrm{id}_{\Delta(X_m)}$.
For any $\mu\in\Delta(\Omega_M)$, define its resolution-$m$ canonical reflection (fiber uniformization) as
\[
\mu^{(m)}:=Q_{M\to m}P_{M\to m}\mu\in\Delta(\Omega_M).
\]
When $m=M$, the above construction degenerates to the $P_M,Q_M$ and $\mu^{e}$ of Section \ref{subsec:pom-information-ledger} (differing only in notation).
\end{definition}

\begin{theorem}[KL Pythagorean identity: residuals are additive on the resolution tower]\label{thm:pom-kl-pythagoras-tower}
Adopt Definition \ref{def:pom-resolution-reflector} and denote KL divergence in the convention of Theorem \ref{thm:pom-kl-ledger}.
Then for any $M\ge m_2\ge m_1\ge 1$ and any $\mu\in\Delta(\Omega_M)$, the strict identity holds:
\[
\boxed{\;
D_{\mathrm{KL}}(\mu\|\mu^{(m_1)})
=D_{\mathrm{KL}}(\mu\|\mu^{(m_2)})
+D_{\mathrm{KL}}(\mu^{(m_2)}\|\mu^{(m_1)})\; }.
\]
\end{theorem}
\begin{proof}
By the definition of KL divergence,
\[
D_{\mathrm{KL}}(\mu\|\mu^{(m_1)})-D_{\mathrm{KL}}(\mu\|\mu^{(m_2)})
=\sum_{\omega\in\Omega_M}\mu(\omega)\log\frac{\mu^{(m_2)}(\omega)}{\mu^{(m_1)}(\omega)}.
\]
For a fixed $y\in X_{m_2}$, on the fiber $(r^{\mathrm{e}}_{M\to m_2})^{-1}(y)$, the ratio
\(
\mu^{(m_2)}(\omega)/\mu^{(m_1)}(\omega)
\)
is constant (since $\mu^{(m)}$ is uniform on each $r^{\mathrm{e}}_{M\to m}$-fiber).
Grouping by $y$ therefore yields
\[
\sum_{\omega\in\Omega_M}\mu(\omega)\log\frac{\mu^{(m_2)}(\omega)}{\mu^{(m_1)}(\omega)}
=\sum_{y\in X_{m_2}}(P_{M\to m_2}\mu)(y)\,
\log\frac{\mu^{(m_2)}(\omega_y)}{\mu^{(m_1)}(\omega_y)},
\]
where $\omega_y$ is any representative point in that fiber.
On the other hand, from the definition \(\mu^{(m_2)}=Q_{M\to m_2}P_{M\to m_2}\mu\) it follows that
\(
P_{M\to m_2}\mu^{(m_2)}=P_{M\to m_2}\mu
\),
so the coefficient $(P_{M\to m_2}\mu)(y)$ on the right-hand side can be replaced by $(P_{M\to m_2}\mu^{(m_2)})(y)$.
``Unfolding'' the grouping back to $\Omega_M$ then gives
\[
\sum_{\omega\in\Omega_M}\mu(\omega)\log\frac{\mu^{(m_2)}(\omega)}{\mu^{(m_1)}(\omega)}
=\sum_{\omega\in\Omega_M}\mu^{(m_2)}(\omega)\log\frac{\mu^{(m_2)}(\omega)}{\mu^{(m_1)}(\omega)}
=D_{\mathrm{KL}}(\mu^{(m_2)}\|\mu^{(m_1)}).
\]
Rearranging yields the conclusion.
\end{proof}

\begin{corollary}[Cross-resolution information conservation ledger: telescoping decomposition]\label{cor:pom-kl-ledger-telescope}
Under the setting of Theorem \ref{thm:pom-kl-pythagoras-tower}, define the single-step externalization
\[
\Delta_m(\mu):=D_{\mathrm{KL}}(\mu^{(m)}\|\mu^{(m-1)})\qquad(2\le m\le M).
\]
Then for any $M\ge 1$ the strict telescoping decomposition holds:
\[
\boxed{\;
D_{\mathrm{KL}}(\mu\|\mu^{(1)})
=\sum_{m=2}^{M}\Delta_m(\mu)\;+\;D_{\mathrm{KL}}(\mu\|\mu^{(M)})\; }.
\]
\end{corollary}
\begin{proof}
Applying Theorem \ref{thm:pom-kl-pythagoras-tower} successively for $m=2,3,\dots,M$ with the triple $(m_1,m_2)=(m-1,m)$ and summing, all intermediate terms cancel telescopically, yielding the stated identity.
\end{proof}

\begin{remark}[Structural criterion: pseudo-invariants arise from coordinate inconsistency]\label{rem:pom-gauge-criterion}
If an ISA/microcode/kernel uses naive truncation ($\tau$) rather than $P_Z$ (or equivalently $r^{\mathrm{e}}$) to compare objects across resolutions, then its ``regularities/invariants'' may be mere phantoms of coordinate inconsistency: a single high-resolution microstate is misclassified at lower resolution, thereby manufacturing pseudo-prime shadows and pseudo-convergence phenomena. Consequently, cross-resolution auditability requires the fold-aware restriction as a gauge condition.
\end{remark}

\begin{theorem}[Inverse limit naturality]\label{thm:pom-inverse-limit}
Suppose that for each resolution $m$ a stable-level operation $f_m:X_m\to X_m$ is defined and satisfies, for all $m_2\ge m_1$, the truncation projection condition
\[
\pi_{m_2\to m_1}\circ f_{m_2}=f_{m_1}\circ \pi_{m_2\to m_1}.
\]
Then there exists a unique limit operation $f_\infty:X_\infty\to X_\infty$ such that $\pi_m\circ f_\infty=f_m\circ \pi_m$ for all $m$.
\end{theorem}
\begin{proof}
By Theorem \ref{thm:inverse-limit-golden}, $X_\infty\cong\varprojlim X_m$. The compatibility identities are precisely the natural transformation condition on the inverse system, and therefore the unique $f_\infty$ is provided by the universal property of the inverse limit.
\end{proof}

\begin{proposition}[Fiber invariant discrimination]\label{prop:pom-fiber-invariance}
Given $g:\Omega_m\to\mathbb{R}$, the following conditions are equivalent:
\begin{enumerate}
  \item If $\Fold_m(a)=\Fold_m(b)$ then $g(a)=g(b)$;
  \item There exists $\bar g:X_m\to\mathbb{R}$ such that $g=\bar g\circ\Fold_m$;
  \item $g$ (as an observable) belongs to the fold-invariant subalgebra $\mathcal{B}_m$ (Proposition \ref{prop:fold-invariant-subalgebra}).
\end{enumerate}
\end{proposition}
\begin{proof}
$(1)\Rightarrow(2)$: Define $\bar g(x)=g(a)$ for any $a\in\Fold_m^{-1}(x)$; (1) guarantees well-definedness. $(2)\Rightarrow(3)$: $g$ is constant on each fiber, hence corresponds to a coordinate representation of $\mathcal{B}_m$. $(3)\Rightarrow(1)$: Elements of $\mathcal{B}_m$ are fiber-constant (Proposition \ref{prop:fold-invariant-subalgebra}).
\end{proof}

\begin{remark}[Time--space duality interface and prime shadows]\label{rem:pom-time-space}
The strict commuting diagram of Proposition \ref{prop:time-space-commuting-diagram} aligns the time-side $\Tr(T^n)$ with the space-side primitive orbit spectrum, thereby providing an auditable interface for the ``time--space orthogonal duality.'' On the other hand, the primitive orbit spectrum of a finite kernel is not identical to the integer primes (see Remark \ref{rem:prime-shadow-vs-primes}): the former is the intrinsic primality of the kernel dynamics, while the latter requires an additional family of prime-index operators. The prime register fiber $\mathcal{P}$ is precisely the ontological supplement for this additional indexing structure.
\end{remark}

\begin{theorem}[Chebotarev lifting on the profinite axis]\label{thm:pom-profinite-axis}
Let $\{G_m,\pi_{m_2\to m_1}\}$ be a projective system of finite groups, $G_\infty:=\varprojlim_m G_m$. For each $m$, consider the corresponding $G_m$-extension kernel and denote
$$
\lambda_m:=\rho(B_{\mathbf 1,m}),\qquad
\Lambda_m:=\max_{\rho\neq\mathbf 1}\rho(B_{\rho,m}).
$$
Let \(U\subset G_\infty\) be a cylinder set, i.e., there exists some $m$ and a conjugation-invariant subset $A\subseteq G_m$ such that
\(
U=\pi_m^{-1}(A)
\).
Then the corresponding primitive orbit count satisfies
$$
B_{n,U}
=\mu_{\mathrm{Haar}}(U)\,\frac{\lambda_m^n}{n}
 +O\!\left(\frac{\Lambda_m^n}{n}\right).
$$
where \(\mu_{\mathrm{Haar}}(U)=|A|/|G_m|\).
If the kernel-version GRH condition $\Lambda_m\le \lambda_m^{1/2}$ is further satisfied (for that level $m$), then the error improves to square-root level:
$$
B_{n,U}
=\mu_{\mathrm{Haar}}(U)\,\frac{\lambda_m^n}{n}
 +O\!\left(\frac{\lambda_m^{n/2}}{n}\right).
$$
\end{theorem}
\begin{proof}
Since $U$ is a cylinder, there exists $m$ and a conjugation-invariant subset $A\subseteq G_m$ such that $U=\pi_m^{-1}(A)$. Let $B_{n,c}$ denote the primitive orbit count falling in the conjugacy class $c\subseteq G_m$; then
$$
B_{n,U}=\sum_{c\subseteq A} B_{n,c}.
$$
For each $c$, the finite-level Chebotarev-type estimate (Theorem \ref{thm:kernel-chebotarev-exp}) gives
$
B_{n,c}=\frac{|c|}{|G_m|}\frac{\lambda_m^n}{n}+O(\Lambda_m^n/n)
$.
Summing over all classes yields the leading coefficient
$
\sum_{c\subseteq A}|c|/|G_m|=|A|/|G_m|=\mu_{\mathrm{Haar}}(U)
$, and since the number of conjugacy classes in $A$ is finite, the error remains $O(\Lambda_m^n/n)$. If $\Lambda_m\le \lambda_m^{1/2}$, the square-root level error follows immediately.
\end{proof}

\begin{corollary}[Cylinder-level effective uniformization length (logarithmic budget for relative error)]\label{cor:pom-profinite-cylinder-effective-n}
Adopt the notation of Theorem \ref{thm:pom-profinite-axis}. Let \(U=\pi_m^{-1}(A)\subseteq G_\infty\) be a cylinder at level \(m\), and assume \(A\neq\varnothing\).
Then for primitive orbit counts of length \(n\), the relative error bound is
\[
\frac{\bigl|B_{n,U}-\mu_{\mathrm{Haar}}(U)\,\lambda_m^n/n\bigr|}{\mu_{\mathrm{Haar}}(U)\,\lambda_m^n/n}
\;=\;
O\!\left(\frac{|G_m|}{|A|}\left(\frac{\Lambda_m}{\lambda_m}\right)^n\right)
\;\le\;
O\!\left(|G_m|\left(\frac{\Lambda_m}{\lambda_m}\right)^n\right).
\]
In particular, if there exists \(\varepsilon>0\) such that the level satisfies the uniform spectral gap condition
\[
\Lambda_m\ \le\ \lambda_m^{1/2-\varepsilon},
\]
then for any target relative error threshold \(\delta\in(0,1)\), it suffices that
\[
\boxed{\;
n\ \gtrsim\ \frac{\log(|G_m|/\delta)}{(1/2+\varepsilon)\log \lambda_m}\; }
\]
to ensure that the Chebotarev leading term on this cylinder achieves relative error of scale \(O(\delta)\).
\end{corollary}

\begin{proof}
By Theorem \ref{thm:pom-profinite-axis},
\(
|B_{n,U}-\mu_{\mathrm{Haar}}(U)\lambda_m^n/n|\le C\,\Lambda_m^n/n
\)
(where the constant \(C\) depends only on the finite kernel at that level and the number of conjugacy classes in \(A\)).
Dividing both sides by the leading term \(\mu_{\mathrm{Haar}}(U)\lambda_m^n/n=(|A|/|G_m|)\lambda_m^n/n\) gives
\[
\frac{|B_{n,U}-\mu_{\mathrm{Haar}}(U)\lambda_m^n/n|}{\mu_{\mathrm{Haar}}(U)\lambda_m^n/n}
\le
\frac{C|G_m|}{|A|}\left(\frac{\Lambda_m}{\lambda_m}\right)^n.
\]
If further \(\Lambda_m\le \lambda_m^{1/2-\varepsilon}\), then \((\Lambda_m/\lambda_m)^n\le \lambda_m^{-(1/2+\varepsilon)n}\).
Setting the right-hand side \(\le \delta\) and solving for \(n\) yields the stated logarithmic budget.
\end{proof}

\begin{corollary}[Inverse: upper bound on supportable cylinder/register capacity given time \(n\)]\label{cor:pom-profinite-cylinder-capacity}
Adopt the notation and uniform spectral gap assumption of Corollary \ref{cor:pom-profinite-cylinder-effective-n}:
\(
\Lambda_m\le \lambda_m^{1/2-\varepsilon}
\).
Fix a target relative error threshold \(\delta\in(0,1)\) and a length budget \(n\).
Then to compress the Chebotarev relative error to \(O(\delta)\) on any nonempty cylinder \(U=\pi_m^{-1}(A)\) (\(A\neq\varnothing\)) at level \(m\), a sufficient capacity condition (in the \(\lesssim\) convention) is
\[
\boxed{\;
\log |G_m|\ \lesssim\ (1/2+\varepsilon)\,n\log\lambda_m\ -\ \log(1/\delta)\; }.
\]
In other words: under the uniform spectral gap convention, the ``supportable register capacity'' \(\log|G_m|\) on the profinite axis is linearly bounded above by the ``time step count'' \(n\).
\end{corollary}
\begin{proof}
From the derivation in the proof of Corollary \ref{cor:pom-profinite-cylinder-effective-n} (see the explicit bound therein),
\(
\mathrm{Err}\lesssim |G_m|\,\lambda_m^{-(1/2+\varepsilon)n}
\).
Setting the right-hand side \(\le \delta\) and taking logarithms yields the stated capacity upper bound (absorbing the implicit constant into \(\lesssim\)).
\end{proof}

\begin{theorem}[Sync-subtracted Chebotarev capacity law]\label{thm:pom-sync-subtracted-chebotarev-capacity}
Fix a level \(m\) and two confidence parameters \(\delta_{\mathrm{NULL}},\delta_{\mathrm{Cheb}}\in(0,1)\).
Adopt the right-resolving deterministic representation of the synchronization kernel and the synchronization time \(T_m\) (Definition \ref{def:Ym-sync-time}), and under the semantic convention that ``address precedes conceptual value / no address until synchronized,'' interpret \(T_m\) as the vanishing moment of the statistical type \(\mathsf{NULL}\) (Remark \ref{rem:Ym-stat-null-unsync}).
\par
Define the synchronization deduction term
\[
n_{\mathrm{sync}}(m,\delta_{\mathrm{NULL}})
:=
\begin{cases}
D_{\mathrm{sync}}(m), & \text{if }Y_m^{\mathrm{amb}}=\varnothing,\\[6pt]
\left\lceil \dfrac{1}{\varepsilon_m}\log\dfrac{C_m}{\delta_{\mathrm{NULL}}}\right\rceil,
& \text{if }Y_m^{\mathrm{amb}}\neq\varnothing,
\end{cases}
\]
where \(D_{\mathrm{sync}}(m)\) is given by Corollary \ref{cor:Ym-sync-delay} (or more generally by the DAG longest path audit convention of Remark \ref{rem:Ym-amb-empty});
the constants \(C_m,\varepsilon_m\) in the second line come from the exponential tail budget of Theorem \ref{thm:Ym-sync-tail}.
Let the available depth be
\[
n_{\mathrm{eff}}:=\bigl(n-n_{\mathrm{sync}}(m,\delta_{\mathrm{NULL}})\bigr)_+,
\qquad (x)_+:=\max\{x,0\}.
\]
\par
If the level satisfies the uniform spectral gap assumption
\(
\Lambda_m\le \lambda_m^{1/2-\varepsilon_{\mathrm{gap}}}
\)
(for some \(\varepsilon_{\mathrm{gap}}>0\)), then any attempt to simultaneously achieve, within a total scan budget of \(n\):
\begin{enumerate}
  \item \textbf{Readout definedness:} \(\nu_m(T_m>n_{\mathrm{sync}})\le \delta_{\mathrm{NULL}}\) (so that the address is prefix-uniquified with probability \(\ge 1-\delta_{\mathrm{NULL}}\)), and
  \item \textbf{Readout statistical clarity:} compressing the Chebotarev relative error to \(\le O(\delta_{\mathrm{Cheb}})\) for any nonempty cylinder \(U=\pi_m^{-1}(A)\),
\end{enumerate}
must use \(n_{\mathrm{eff}}\) as the ``available depth'' entering the Chebotarev budget. Equivalently, the capacity condition of Corollary \ref{cor:pom-profinite-cylinder-capacity} undergoes a synchronization deduction correction:
\[
\boxed{\;
\log |G_m|
\ \lesssim\
(1/2+\varepsilon_{\mathrm{gap}})\,n_{\mathrm{eff}}\log\lambda_m
\ -\
\log(1/\delta_{\mathrm{Cheb}})
\;}
\]
(the implicit constant convention is consistent with Corollary \ref{cor:pom-profinite-cylinder-capacity}).
\par
\noindent\textbf{Auditability note (as part of the conclusion).}
\begin{enumerate}
  \item The synchronization deduction term \(n_{\mathrm{sync}}\) can be directly audited on the finite determinized graph: the ambiguity subgraph contains a directed cycle if and only if \(Y_m^{\mathrm{amb}}\neq\varnothing\); if acyclic, then \(D_{\mathrm{sync}}\) is obtained by longest-path DP (Remark \ref{rem:Ym-amb-empty} and Table \ref{tab:phi_m_sync_theory}).
  \item The synchronization delay is a non-spatializable time invariant: it is orthogonal to any extension of residue/prime-register and other orthogonal recording axes, and cannot be eliminated by merely increasing modular axis capacity; therefore it must enter any capacity budget as an explicit time deduction term (see Remark \ref{rem:pom-time-defect-dichotomy}).
\end{enumerate}
\end{theorem}

\begin{corollary}[Sharp linear deduction under Zeckendorf folding]\label{cor:pom-sync-subtracted-chebotarev-zeckendorf}
In the Zeckendorf folding model of this paper, when \(m\ge 3\) we have \(Y_m^{\mathrm{amb}}=\varnothing\) and \(D_{\mathrm{sync}}(m)=m-1\) (Corollary \ref{cor:Ym-sync-delay} and Remark \ref{rem:Ym-amb-empty}).
Therefore for any \(\delta_{\mathrm{NULL}},\delta_{\mathrm{Cheb}}\in(0,1)\),
\[
n_{\mathrm{eff}}=(n-m+1)_+,
\]
and the capacity condition of Theorem \ref{thm:pom-sync-subtracted-chebotarev-capacity} degenerates to the sharp linear deduction form
\[
\log |G_m|
\ \lesssim\
(1/2+\varepsilon_{\mathrm{gap}})\,(n-m+1)_+\log\lambda_m
\ -\
\log(1/\delta_{\mathrm{Cheb}}).
\]
\end{corollary}

\begin{corollary}[Two-parameter budget in the weakly reversible case]\label{cor:pom-sync-subtracted-chebotarev-weak}
If \(Y_m^{\mathrm{amb}}\neq\varnothing\), then to make the statistical type \(\mathsf{NULL}\) vanish under the maximal entropy measure with probability \(\le \delta_{\mathrm{NULL}}\), one must first pay the synchronization budget
\[
n_{\mathrm{sync}}(m,\delta_{\mathrm{NULL}})
=\left\lceil \dfrac{1}{\varepsilon_m}\log\dfrac{C_m}{\delta_{\mathrm{NULL}}}\right\rceil,
\]
after which a uniform Chebotarev statistical budget becomes meaningful; and the time-side capacity condition couples two exponential rates via the deduction \(n_{\mathrm{eff}}=(n-n_{\mathrm{sync}})_+\): \(\varepsilon_m\) controls the synchronization tail rate for definitional clarity, while \(\lambda_m\) and \(\varepsilon_{\mathrm{gap}}\) control the spectral rate for statistical clarity (Theorem \ref{thm:pom-sync-subtracted-chebotarev-capacity}).
\end{corollary}

\subsubsection{Lift Operator on \texorpdfstring{$L^2(\widehat{\ZZ})$}{L2(Zhat)}: Character Direct-Sum Decomposition and an Operator Framework ``Unifying All Moduli''}\label{subsec:profinite-operator-l2}
Theorem \ref{thm:pom-profinite-axis} provides the Chebotarev counting and error convention on a per-finite-level \(G_m\) basis.
This subsection gives a compression closer to the Hilbert--P{\'o}lya semantics: unifying all modular twist channels into a single lift operator acting on \(L^2(G_\infty)\), and strictly diagonalizing it on the Pontryagin dual side into a family of twist matrices.

\paragraph{Specialization to the prime register axis: \texorpdfstring{$G_\infty=\widehat{\ZZ}$}{G=Zhat}}
Take
\[
G_\infty=\widehat{\ZZ}=\varprojlim_m \ZZ/m\ZZ,
\]
which is a compact abelian group whose Pontryagin dual is the discrete group
\[
\widehat{G_\infty}\cong \QQ/\ZZ,
\]
possessing a standard orthonormal character basis \(\{\chi_\alpha\}_{\alpha\in\QQ/\ZZ}\subset L^2(\widehat{\ZZ})\); see the main text discussion on character bases and duality, as well as the DFT convention for congruence axes in the appendix.

\paragraph{Lift operator}
Let the state set of the kernel be the finite set \(V\), and denote the weighted directed edge set by \(E\subset V\times V\).
To each edge \(e:i\to j\) assign a weight \(w(e)\in\CC\) (given, for example, by the output potential / collision potential / temperature gate), and assign a ``register increment'' label
\(\gamma(e)\in \widehat{\ZZ}\)
(at a finite level \(\ZZ/m\ZZ\) this is simply the congruence update step; at the inverse limit it is uniquely determined by compatibility).
Define the Hilbert space
\[
\mathcal{H}:=\ell^2(V)\otimes L^2(\widehat{\ZZ}),
\]
and define the lift operator \(T:\mathcal{H}\to\mathcal{H}\) by
\[
\boxed{\ (Tf)_j(x):=\sum_{e:i\to j} w(e)\,f_i\!\bigl(x-\gamma(e)\bigr)\qquad(x\in\widehat{\ZZ})\ }.
\]
This is precisely the standard ``finite state \(\times\) compact group translation'' construction: it couples the finite kernel transition with the register translation into a single operator.

\paragraph{Character decomposition: simultaneous diagonalization of all twist channels}
Since \(\widehat{\ZZ}\) is a compact abelian group, the Fourier transform gives a unitary isomorphism
\[
L^2(\widehat{\ZZ})\ \cong\ \ell^2(\QQ/\ZZ),
\]
and diagonalizes the translation action into character multipliers. Thus \(T\) decomposes on the Fourier side into a \textbf{strict direct sum} (viewable as a direct integral with respect to counting measure):
\[
\boxed{\ T\ \simeq\ \bigoplus_{\alpha\in\QQ/\ZZ} M_{\chi_\alpha}\ },
\]
where each fiber twist matrix \(M_{\chi}\in\CC^{V\times V}\) has entries
\[
\boxed{\ (M_\chi)_{ij}:=\sum_{e:i\to j} w(e)\,\chi\!\bigl(\gamma(e)\bigr)\ }.
\]
Therefore ``twisting over all moduli / all characters'' is unified as the character-component spectrum of a single operator \(T\).
At a finite level \(G_m=\ZZ/m\ZZ\), the above formula degenerates to block diagonalization under the discrete Fourier transform (DFT), which is precisely the linear algebra mechanism behind Theorem \ref{thm:pom-profinite-axis} and the finite-level Chebotarev estimates.

\begin{remark}[Grand Unified Lift: incorporating channel/boundary tower into lift operator fibers]\label{rem:pom-grand-unified-lift}
The \(T\) above unifies all twist channels on the ``congruence/modular axis.''
If one further wishes to incorporate the \((\varphi,\pi,e)\) three projection channels and the multi-resolution boundary tower into the same operator framework,
one can introduce an additional compact group (or finite group) fiber \(G_{\mathrm{chan}}\) to carry these discrete labels, and define the extended fiber
\[
G_\infty^{\mathrm{GUT}}
:=\widehat{\ZZ}\times G_{\mathrm{chan}},
\qquad
\mathcal{H}_{\mathrm{GUT}}:=\ell^2(V)\otimes L^2\!\left(G_\infty^{\mathrm{GUT}}\right).
\]
Let each edge \(e:i\to j\) carry, in addition to \(\gamma(e)\in\widehat{\ZZ}\), a channel increment \(\eta(e)\in G_{\mathrm{chan}}\).
Then the corresponding lift operator can be written as
\[
(T_{\mathrm{GUT}}f)_j(x,h)
:=\sum_{e:i\to j} w(e)\,f_i\!\bigl(x-\gamma(e),\ h\cdot \eta(e)\bigr),
\qquad (x,h)\in \widehat{\ZZ}\times G_{\mathrm{chan}}.
\]
On the dual side, the character decomposition of \(\widehat{\ZZ}\) and the Peter--Weyl decomposition of \(G_{\mathrm{chan}}\) simultaneously block-diagonalize it into finite-dimensional twist blocks labeled by \((\chi,\varrho)\):
\[
T_{\mathrm{GUT}}\ \simeq\ \bigoplus_{\chi\in\widehat{\ZZ}}\ \bigoplus_{\varrho\in\widehat{G_{\mathrm{chan}}}}\ M_{\chi,\varrho},
\qquad
(M_{\chi,\varrho})_{ij}:=\sum_{e:i\to j} w(e)\,\chi(\gamma(e))\,\varrho(\eta(e)).
\]
Thus ``grand unification'' can be formulated entirely at the level of finite-dimensional, enumerable spectral data: for a family of channels \((\chi,\varrho)\) activated by observation gates, the spectral gap coupling \(g_{\chi,\varrho}(u)\) achieves alignment at a certain unification point (see the minimax criterion of Section \ref{subsec:group-unification-spectral-alignment}).
\end{remark}

\begin{remark}[Adelic anomaly functional: a falsifiable unified criterion for local--global anomaly cancellation]\label{rem:pom-adelic-anom-functional}
The spectral gap alignment (Section \ref{subsec:group-unification-spectral-alignment}) provides a unified convention for exponential-level slow modes;
if one also wishes to incorporate the \emph{constant-level residuals} from the non-commutativity of ``\(E_m\)--\(\mathrm{LIFT}\)'' into the same computable framework,
one can, on top of the anomaly signature \(\mathsf{Anom}_G(K;\theta)\) of Definition \ref{def:pom-anom-signature}, introduce for each finite level (e.g., the prime modular axis \(G_p=\ZZ/p\ZZ\)) a scalar local anomaly strength.
An optimization/gradient-friendly choice is
\[
A_p(\theta):=\bigl\|\mathsf{Anom}_{G_p}(K;\theta)\bigr\|_{\ell^2}^2
=\sum_{\chi\in\widehat{G_p}\setminus\{\chi_0\}}\bigl|\log\mathfrak{M}_\chi(\theta)\bigr|^2,
\]
where \(\chi_0\) is the trivial character (alternatively, the \(\ell^\infty\) norm \(\max_{\chi\ne\chi_0}|\log\mathfrak{M}_\chi(\theta)|\) can serve as an audit convention).
Under the profinite convention \(G_\infty=\widehat{\ZZ}\), taking a prime truncation \(P\), define the global anomaly functional
\[
\mathcal{A}_{\le P}(\theta):=\sum_{p\le P} w_p\,A_p(\theta),
\]
where the weights \(w_p\) can be taken as \(\log p\), \(p^{-2}\), etc., to reflect scale normalization or channel budgeting.
Then ``cross-modular consistent lifting / anomaly cancellation'' can be formulated as a fully falsifiable, script-auditable unified criterion:
if there exists \(\theta^\star\) such that \(\mathcal{A}_{\le P}(\theta)\) achieves a stable minimum near \(\theta^\star\) (or satisfies \(\nabla_\theta\mathcal{A}_{\le P}(\theta^\star)\approx 0\)),
and as the truncation \(P\) grows this minimum point and value remain numerically stable (robust to truncation depth/sampling strategy),
then \(\theta^\star\) can be regarded as a ``local--global adelic flat point'' (i.e., global cancellation of constant-level curvature/anomaly).
Since each \(A_p(\theta)\) can be computed via the \(\zeta\)/finite-part interface of finite-dimensional twist blocks \(M_\chi\), the above functional and stability conditions can be directly translated into reproducible optimization and regression certificates.
\end{remark}

\paragraph{Normalized spectrum and GRH semantics}
Denote the Perron root of the trivial channel by \(\lambda:=\rho(M_{\mathbf{1}})\), and for any character \(\chi\) define the normalized fiber operator
\[
U_\chi:=\lambda^{-1/2}M_\chi.
\]
Let \(w_{\chi,1},\dots,w_{\chi,|V|}\) be the eigenvalues of \(U_\chi\); then the ``critical circle'' condition is \(|w_{\chi,i}|=1\).
By the spectrum--\(s\)-plane dictionary of \S\ref{subsec:zeta-online-kernel} (Proposition \ref{prop:s-plane-spectrum-poles}) and the toy-RH discussion in Appendix \ref{app:real-input-40-finite-rh}, this is equivalent to confining the corresponding poles to \(\Re(s)=\tfrac12\).
Therefore an operational ``spectral unitarization index'' for the kernel-version GRH can be taken as
\[
\mathcal{E}(\chi):=\max_{1\le i\le |V|}\bigl|\log|w_{\chi,i}|\bigr|,
\]
and the ``GRH-type assertion'' is understood as: within the character domain/scale domain of interest, \(\mathcal{E}(\chi)\) is sufficiently small (or tends to zero in the limit).

\paragraph{Double-scale barrier and diffusion convention}
When the character family is extended to \(\QQ/\ZZ\) (including arbitrarily high-order small-angle characters), the ``worst-case twisted spectral radius ratio'' generally cannot maintain a uniform bound away from \(1\) across all moduli.
In the synchronization-kernel weighted case, the appendix has already derived the near-coboundary small-angle quadratic law and the diffusion threshold \(n\asymp m^2\) (Corollary \ref{cor:sync-rho-gap-m8-diffusive}).
This phenomenon suggests that under the profinite inverse limit semantics, a more natural GRH/Chebotarev formulation should employ a double-limit scale in both length \(n\) and modulus \(m\) (cf.\ the abstract formulation in Conjecture \ref{conj:scaled-kernel-grh-diffusive}).
